\scanpage{491}
\mspage{249}

\uline{35} \uline{L'injectivité de} $\Gamma_{\mathbb Q}\longrightarrow\mathrm{Autext}_{\mathrm{lac}}(\widehat{\Pi}_{0,3})$.

\margin{démontré modulo le lemme 2 plus bas ---}

\uline{Théorème 1} Soit $\Gamma=\mathrm{Gal}(\overline{\mathbb Q}_0/\mathbb Q)$, où $\overline{\mathbb Q}_0$ est la clôture algébrique de $\mathbb Q$ dans $\mathbb C$, considérons l'homomorphisme canonique
\[
(1)\qquad \unsure{\Gamma}\xrightarrow{\Theta}\widehat{\widehat{\unsure{\mathcal T}}}_{0,3}=\mathrm{Autext}_{\mathrm{lac}}(\widehat{\Pi}_{0,3}).
\]
Cet homomorphisme est \uline{injectif}.

Démonstration.

\uline{Lemme 1} Soit $x\in\mathbb Q$, $x\ne0,1$ i.e. $x\in U_{0,3}(\mathbb Q)$, choisissons un \deleted{lacet} chemin sur $\mathbb P^1(\mathbb C)$ de $P=-\unsure{\overline{j}|\overline{J}}$ (point base pour définir $\Pi_{03}=\pi_1(U_{0,3}(\unsure{\mathbb C}),P)$ --- N.B. $U_{03}=\mathbb P^1_{\mathbb Q}\setminus\{0,1,\infty\}$) d'où un isomorphisme $\pi_1(U_{0,3\,\overline{\mathbb Q}_0},x)\simeq\widehat{\Pi}_{0,3}$, et par transport de structure une action \inserted{effective} de $\Gamma$ sur $\widehat{\Pi}_{0,3}$ (pas seulement extérieure). Alors $K\unsure{=}\mathrm{Ker}\,\Theta$ opère trivialement sur $\widehat{\Pi}_{0,3}$.

\inserted{Dém. Il suffit de voir que l'opération \uline{extérieure} de $K$ sur $\pi_1(\overline V)$, où $V=U_{0,3}\setminus\{x\}$, est triviale.}
D'après le résultat de Белый, il existe un morphisme (\uline{défini sur} $\mathbb Q$, ceci est essentiel)
\[
\mathbb P^1_{\mathbb Q}\longrightarrow\mathbb P^1_{\mathbb Q},
\]
\deleted{\illegible{1}}\unsure{étale} au dessus de $U_{0,3}=\mathbb P^1_{\mathbb Q}\setminus\{0,1,\infty\}$, et tel que $x\longmapsto0$. Soit $U'$ le revêtement étale de $U_{03}=U$ induit, \deleted{\illegible{2}} $\pi'=\pi_1(\overline{U'})$ son groupe fondamental géo-
\margin{490}

\scanpage{492}

métrique, sur lequel $\Gamma$ donc $K\subset\Gamma$ opère extérieurement. Alors $\pi(\overline V)$ est un quotient de $\pi'$, avec respect des opérations extérieures de $\Gamma$, et il suffit de prouver que $K$ opère trivialement sur $\pi'$. Mais $\pi'$ s'identifie à un sous-groupe ouvert de $\widehat{\Pi}_{0,3}=\pi$ (avec respect des opérations extérieures de $K$, sur lequel l'opération extérieure de $K$ est triviale. Il s'ensuit que l'opération \inserted{ext.} de $K$ sur $\pi'$ se fait à travers un groupe quotient \uline{fini}.

\margin{\diagram{\begin{tikzcd}[ampersand replacement=\&]
1 \arrow[r] \& \pi' \arrow[r] \arrow[d] \& \mathcal E'_K \arrow[r] \arrow[d] \& K \arrow[r,dash] \arrow[d,equal] \& 1 \\
1 \arrow[r,dash] \& \pi \arrow[r] \& \mathcal E_K \arrow[r,dash] \& K \arrow[r,dash] \& 1
\end{tikzcd}}}

[En effet, l'action extérieure de $K$ se fait à travers l'action effective du groupe extension $\mathcal E_K$, laquelle par construction de $K$ se fait par un homomorphisme $\mathcal E_K\longrightarrow\pi/\unsure{Z|\mathcal Z}$ ($Z=\mathrm{Centre}(\pi)$), \deleted{\illegible{1}} l'action de $\mathcal E'_K$ \unsure{induit|induite} se fait par le composé $\mathcal E'_K\longrightarrow\mathcal E_K\longrightarrow\pi/Z$, dont l'image est dans $N/Z$, où $N$ est le normalisateur de $\pi'$ dans $\pi$. Donc l'image de $K$ dans $\mathrm{Autext}(\pi')$ est contenue dans celle de $N/\pi'\longrightarrow\mathrm{Autext}(\pi')$, or $N/\pi'$ est fini ]
\margin{491}

\scanpage{493}
\mspage{250}

Repassant à $\pi_1(\overline V)$, on voit donc que l'image de $K$ dans $\mathrm{Autext}\,\pi_1(\overline V)$ est finie, donc l'image de $K$ dans $\mathrm{Aut}\,\widehat{\Pi}_{03}$ est finie. Elle est d'ailleurs formée d'automorphismes intérieurs, donc le Lemme 1 sera conséquence du

\uline{Lemme 2} Tout automorphisme intérieur d'ordre fini de $\widehat{\Pi}_{0,3}$ (groupe profini libre à deux générateurs) est trivial. De façon plus précise : $\widehat{\Pi}_{0,3}$ a un centre réduit à $\{1\}$, et tout élément de $\widehat{\Pi}_{0,3}$ d'ordre fini est réduit à 1.

Ceci est un énoncé d'algèbre profinie pure, que je reporte pour plus tard, pour en terminer avec la partie ``géométrique'' de la démonstration du théorème.

\uline{Lemme} 3 Pour \unsure{tout|toute} ouvert \inserted{non vide} $\deleted{U}V\subset\deleted{\unsure{U_{0,3}}}\inserted{\mathbb P^1_{\mathbb Q}}$, \inserted{considérant} l'action extérieure de $\Gamma$ sur $\pi_1(\overline{\deleted{U}V})$, la restriction de celle-ci à $K$ est triviale.

Quitte à passer à un ouvert plus
\margin{492}

\scanpage{494}

petit, OPS par Белый que $\deleted{U}V$ est un revêtement étale de $U_{03}$ --- et le raisonnement précédent montre alors que l'action extérieure de $K$ se fait via un groupe \inserted{quotient} fini --- mais cela n'est pas suffisant pour notre propos, et n'implique pas par lui-même que cette action soit triviale.

\unsure{(}\deleted{Mais choisis}\inserted{D'ailleurs, au point où j'en suis,} on aurait pu remplacer $U$ par n'importe quelle courbe algébrique qu.proj. lisse géométriquement connexe --- pour trouver que $K$ opère sur le groupe extérieur $\pi_1(\overline{\deleted{U}V})$ via un groupe fini. Mais ici l'hypothèse $\deleted{U}V\subset\mathbb P^1_{\mathbb Q}$ implique qu'il existe $y\in V(\mathbb Q)$, soit $x\in U(\mathbb Q)$ son image dans $U=U_{0,3}$.

Prenant \deleted{\unsure{les}} $y,x$ comme points-bases pour les groupes fondamentaux, on trouve maintenant un homomorphisme effectif de groupes fondamentaux
\margin{493}

\scanpage{495}
\mspage{251}

\[
\pi'=\pi_1(\overline V,y)\hookrightarrow
\pi\;(\simeq\widehat{\Pi}_{0,3})=\pi_1(\overline U,x),
\]
compatible avec une action \uline{effective} de $\Gamma$ sur ces groupes. Par le lemme 1
\deleted{cette}\inserted{l'}action induite de $K$ sur $\pi=\pi_1(\overline U,x)$ est triviale,
donc aussi son action sur le sous-groupe $\pi'$. A fortiori l'action extérieure est triviale, cqfd.

On peut maintenant prouver le

\uline{Lemme 4} $K=(\unsure{1})$ (i.e. le théorème !)

En effet, \inserted{sous les conditions du lemme 2,} l'action de
\deleted{\unsure{$\Gamma$}}\inserted{$K$} sur l'ensemble
$I=S(\overline{\mathbb Q}_0)$, où $S=\mathbb P^1_{\mathbb Q}\setminus V$,
est déduite de l'action extérieure de
\deleted{\unsure{$\Gamma$}}\inserted{$K$} sur $\pi_1(\overline V)$,
\inserted{si $\operatorname{card} I\geq 2$,} comme l'action sur les classes de conjugaison de
``\unsure{ss-groupes|s-groupes} lacets''. Comme $K$ opère trivialement sur le \unsure{groupe} extérieur,
il opère trivialement sur $I$. Ceci \unsure{étant} vrai pour tt $V$, on voit que l'action de $K$
sur $\mathbb P^1(\overline{\mathbb Q}_0)$
\inserted{$=\overline{\mathbb Q}_0\cup\{\infty\}$}
est triviale, i.e. son action sur $\overline{\mathbb Q}_0$ l'est,
donc $K=(1)$, cqfd.

\margin{Ouf !}
\margin{494}

\scanpage{496}

Il reste à reporter la démonstration du lemme 2, que je vais reformuler sous une forme plus générale :

\uline{Théorème 2} Soit $\pi$ un groupe profini libre sur un ensemble \inserted{fini $I$} de cardinal $\geq 2$.

Alors :

a) Le centre de $\pi$ est \deleted{trivial} égal à $(1)$.

b) Il n'y a pas \inserted{dans $\pi$} d'élément d'ordre fini autre que $1$.

Finalement je cale sur a) -- tout comme je ne vois pas pourquoi le centre de $\Gamma$
(et de tt sous-groupe ouvert) doit être réduit à $(1)$. Consulter Deligne à ce sujet !

Pour b), voici un expédient.

Soit $K$ un corps alg. clos. de car. $0$, et considérons le corps des fonctions
$L=K(x)$ de \deleted{\unsure{$U_{0,3\,K}$|$(U_{0,3})_K$}}
\inserted{$U$}$=\mathbb P^1_K\setminus$\deleted{\unsure{$I$|$\Sigma$}}
\inserted{($n+1$ points)}
\deleted{\unsure{($I$} \illegible{1} plongé dans $\mathbb P^1(K)$)}
($n=\operatorname{card} I$). Alors $\pi\simeq\pi_1(U)$, et c'est un quotient de
$\mathcal E_L=\mathrm{Gal}(\overline L/L)$.
Comme $\pi$ est libre, $\mathcal E_L\longrightarrow\pi$ se \unsure{relève} en
$\pi\longrightarrow\mathcal E_L$, et il suffit de voir

\margin{495}

\scanpage{497}
\mspage{252}

que $\mathcal E_L$ n'a pas d'élément d'ordre fini $\ne 1$.

Mais d'après Artin, il n'y a pas d'\unsure{autom.} \inserted{\unsure{$\tau$}}
d'ordre fini \inserted{$\ne 1$} d'un corps alg. clos $\overline L$,
sauf dans le cas où $\overline L$ est extension quadratique d'un
\deleted{\illegible{1}} corps \deleted{\unsure{réel}}\inserted{ordonné} maximal
\inserted{$R$}, \deleted{\unsure{fixé|fini}} qui soit le corps des invariants de
$\tau$ (donc $\tau$ d'ordre $2$ \unsure{exactement}).

\deleted{\illegible{1}} Mais en l'\unsure{occurrence|occurence} on aurait $R\supset L$,
\inserted{et $L$}$\supset K$, or dans $K$ l'élément $-1$ est un carré,
donc $R$ ne peut être ordonné; absurde !

\uline{Corollaire 1 du \unsure{Th.} 1} Soit \unsure{$X$|$U$} une courbe algébrique
(lisse, \unsure{géom\textsuperscript{t}.} connexe, quasi-projective) sur $k$ extension finie de
$\mathbb Q$, \deleted{Alors \unsure{il|A}} (Je présume que l'action extérieure de
$\mathcal E_k^{\overline k}$ sur $\pi_1(\unsure{X|U}_{\overline k})$ est fidèle
si $U$ \unsure{anabélien} -- à défaut de pouvoir le prouver, j'énonce :)
Alors il existe une partie ouverte non vide $V$ de \unsure{$X$|$U$} telle que
l'action \inserted{extérieure} de $\mathcal E_k^{\overline k}$ sur
$\pi_1(V_{\overline k})$ soit fidèle.

\margin{496}

\scanpage{498}

\uline{Dém}. D'après Белый, on sait qu'on peut trouver $V\subset\unsure{X}$, qui soit un revêtement étale de $U_{\deleted{\unsure{03}}}=(U_{03})_k$, je dis que ce $V$-là convient. Soit donc $K$ le noyau de l'\deleted{représentation}\inserted{opération} extérieure de $\Gamma=\mathcal E_k^{\bar k}$ sur $\pi'=\pi_1(V_{\bar k})$. Soit $y$ un pt fermé de $V$, rationnel sur l'extension finie $k'$ de $k$, correspondant à un sous-groupe ouvert $\Gamma'\inserted{=\mathcal E_{k'}^{\bar k}}$ de $\Gamma$, soit $K'=\Gamma'\cap K$. Soit $x$ l'image de $y$ \unsure{dans} $U$, on a
\[
\pi'=\pi_1(V_{\bar k},y)\hookrightarrow\pi=\pi_1(U_{\bar k},x),
\]
et par construction l'action extérieure de $K'$ sur $\pi'$ est triviale, donc $K'$ opère sur $\pi'$ par autom intérieurs. Or on a le

\margin{prouvé seulement modulo vérification de $(2),(4)$ ci-dessous -- pour le \unsure{corollaire} $(2)$ est d'ailleurs suffisant\ldots}
\uline{Lemme} 5 Soit $\pi$ un groupe \inserted{profini} à lacets, $\pi'$ un sous-groupe ouvert, alors tout automorphisme $u$ de $\pi$ qui laisse stable $\pi'$ et induit sur $\pi'$ un autom int. (resp. l'identité) est intérieur (resp. l'identité).

\margin{497}

\scanpage{499}
\mspage{253}

Admettons pour l'instant ce lemme -- il en résulte que les éléments de $K'$, qui opèrent sur $\pi$ en induisant sur $\pi'$ des autom intérieurs, \unsure{induisent|s'induisent} sur $\pi$ lui-même des autom intérieurs, i.e que l'action extérieure de $K'$ sur $\pi$ est triviale. D'après le théorème 1, on sait d'autre part que l'action ext. de $\Gamma'$ sur $\pi$ est fidèle, donc $K'=(1)$. Il en résulte que $K$ est \uline{fini}. Mais par Artin on sait que les seuls sous-groupes finis \inserted{$\ne(1)$} de $\Gamma$ sont ceux engendrés par une "conjugaison \deleted{de $\tau$} complexe" $\tau$, \unsure{correspondante|correspondant} \unsure{à un|aux} sous-corps \deleted{réels}\inserted{\unsure{ordonné|ordonnés}} \unsure{maximal|maximaux} entre $k$ et $\bar k$. Mais pour un tel $\tau$ on a $\chi(\tau)=-1$, donc l'action extérieure de $\tau$ ne peut être triviale -- on gagne. En fait, la démonstration a montré ceci :

\uline{Corollaire 2} Soient $\unsure{k|K}$ \unsure{un} \deleted{extension de}\inserted{corps de car $0$},

\margin{498}

\scanpage{500}

\deleted{\unsure{type} fini de $\mathbb Q$}\inserted{$U$ une courbe} algébrique (lisse etc) sur $k$, \unsure{telles|telle} que $\mathcal E_k^{\bar k}=\Gamma$ opère fidèlement sur $\pi_1(U_{\bar k})$ (opération extérieure).

Alors pour tout revêtement étale géom.\textsuperscript{t} connexe $V$ de $U$, l'opération extérieure de \deleted{\unsure{$\pi$}} $\Gamma$ sur $\pi_1(V_{\bar k})$ est également fidèle.

Reste à prouver le lemme 5. Il suffit de le prouver dans le cas respé -- l'autre s'en déduit aussitôt. \deleted{Mais soit}\inserted{Si $\pi$ a au moins une classe de lacets (cas d'une courbe} \deleted{$(x_i,y_i)_{1\le i\le g}$, $(l_j)_{1\le j\le\nu}$ une base de $\pi$, ---}\inserted{alg. affine i.e. non propre), alors $\pi$ est} \deleted{nous admettrons que si} libre -- et le lemme 5 \deleted{\unsure{se trouve|se prouve}}\inserted{est valable} justement pour de tels groupes. Si $(\unsure{\ell_i|h_i})_{1\le i\le\nu}$ est un syst. de générateurs, soit $L_i$ le sous-groupe fermé engendré par $\unsure{\ell_i|h_i}$, nous admettrons que $\forall n\in\mathbb N^*$
\margin{pas prouvé !}
\[
(2)\qquad L_i=\operatorname{Centr}_{\pi}(\unsure{L_i^n|\ell_i^n})
\]

Ceci posé, si l'autom $u$ de $\pi$ induit l'identité sur $\pi'$, $\exists n\in\mathbb N^*$ tel que $\forall 1\le i\le\unsure{\nu}$
$u(\unsure{\ell_i|h_i}^{n})=\unsure{\ell_i|h_i}^{n}\unsure{{}_r}$, donc $u(\unsure{\ell_i|h_i})$ centralise $\unsure{\ell_i|h_i}^{n}=u(\unsure{\ell_i|h_i})^{n}$,

\margin{499}
